\scanpage{19}
La décomposition envisagée (can, fonct.)
\[
\mathrm{DR}^1(A/S)\simeq\underline{t}_A^\vee+\underline{t}_{A^*}
\]
vaut pour tt préschéma abélien \uline{ordinaire} sur
un $S$ de car. $p>0$, en $=$ temps que les expressions
\[
\underline{t}_A\simeq\unsure{\Phi|\mathscr L|\mathcal L}_A\otimes_{\mathbb F_p}\underline O_S,\qquad
\underline{t}_{A^*}\simeq\Phi_{A^*}\otimes_{\mathbb F_p}\underline O_S
\]
où $\Phi_A,\Phi_{A^*}$ sont des faisceaux \inserted{étales} loc. \deleted{\unsure{constants|const.}}
libres de rang $g$ ($=$ dim rel de $A/S$) d'esp vectoriels
sur $\mathbb F_p$. Si $S$ est le spectre d'un corps parfait,
on a
\[
\begin{aligned}
\Phi_{A^*}=(\underline t_{A^*})^F&=H^1(A,\underline O_A)^F=H^1(A,\mathbb Z/p\mathbb Z)\\
&=\operatorname{Hom}({}_pA\unsure{'|\cdot|{}},\mathbb Z/p\mathbb Z)=\deleted{\unsure{\Phi_A}}{}_pA(k)^\vee
\end{aligned}
\]
\[
\left\{\begin{aligned}
\Phi_{A^*}\deleted{\unsure{k}}(k)&={}_pA(k)^\vee\\
\Phi_A(k)&={}_pA^*(k)^\vee
\end{aligned}\right.
\]
Donc
\[
\left\{\begin{aligned}
\underline t_A^\vee&=\Phi_A\deleted{\unsure{(k)}}^\vee\otimes_{\mathbb F_p}k={}_p A^*(k)\otimes_{\mathbb F_p}k\\
t_{A^*}&={}_pA(k)^\vee\otimes_{\mathbb F_p}k .
\end{aligned}\right.
\]

On a des relations analogues pour $S$ (de car $p>0$)
quelconque : l'hyp. sur $A$ implique que
$\,{}_pA\mathbin{\to}S$ s'écrit de façon unique comme
composé $\,{}_pA\longrightarrow({}_pA)_{\mathrm{\acute et}}\mathbin{-}S$ de deux morphismes
finis loc libres, le premier radiciel surj, le $2^{\mathrm{ième}}$ \uline{étale}.
% archive sheet number: 18
